\subsubsection{Three Functional Groups for the Primary F Ablation}
We group training mechanisms by objective and data dependence before obtaining any additional target scores. Group A is cross-receiver identity contrast (CRIC), including same-TX/different-RX positives, different-TX negatives, joint shallow/intermediate descriptors and their fixed normalization, temperature and ramp. Group B is paired stochastic prediction consistency: the ordinary128, ordinary64 and raw128 symmetric-KL terms switch together. Group C is the entire native selective self-training path: own EMA, grouped confidence selection, neighbor and strong-view agreement, pseudo-label CE and unlabeled entropy. These are two supervised mechanisms and one semi-supervised mechanism; their implementation details do not become separate claimed contributions.

For switches $a,b,c\in\{0,1\}$, the objective is
\begin{equation}
\begin{split}
\mathcal L(a,b,c)={}&\mathcal L_{\rm MVCE}+a\lambda_{\rm RX}\mathcal L_{\rm RX}\\
&+b(0.5\mathcal K_{o128}+0.25\mathcal K_{o64}+0.25\mathcal K_{r128})\\
&+c(0.16\mathcal L_{\rm PL}+0.01\mathcal L_{\rm ent}).
\end{split}
\end{equation}
The common labeled scaffold is $\mathcal L_{\rm MVCE}=0.5\mathcal C_{o128}+0.25\mathcal C_{o64}+0.25\mathcal C_{As64}+0.25\mathcal C_{r128}$, where each $\mathcal C$ averages the original two CE forwards with label smoothing 0.01. Disabling B retains these six labeled forwards and their CE dose; As KL remains zero. Disabling C removes the U loader, EMA and every unlabeled forward/update, rather than evaluating zero-weight unlabeled branches. All 50 epochs still train on L. Backbone, physical branches and the jointly trained default head are fixed. The native T1 As/U recipe is retained for this matched objective study; it does not acquire an additional concatenation loss. The separately registered matched Sup/ordinary-PL and strict architecture controls retain default satellite concatenation.

\begin{table}[t]
\centering\footnotesize
\caption{Primary four-row ablation and the predeclared full factorial. Each cell has three F seeds. All new results are pending.}
\label{tab:three-core-design}
\begin{tabular}{ccp{0.58\columnwidth}}
\toprule
ABC & Main table & Source / purpose\\
\midrule
111 & Yes & Full final T1; existing three weights\\
011 & Yes & Without CRIC; corrected F16\\
101 & Yes & Without paired consistency; corrected F04\\
110 & Yes & Without selective SSL; corrected F20\\
\midrule
000 & Appendix & Matched multiview CE scaffold; new\\
001 & Appendix & Only selective SSL; new\\
010 & Appendix & Only paired consistency; new\\
100 & Appendix & Only CRIC; new\\
\bottomrule
\end{tabular}
\end{table}

The design contains 24 logical models, reusing 12 completed or already scheduled source rows and training only 12 missing combinations. Each new model follows scratch20 plus own10 plus own20, totaling E50/11100 updates; C=0 uses L throughout. Added training is 600 model-epochs/133200 updates. The fixed F0 contract contains L6300, hidden-label U56700 and V27000 from five source RX and three dates, with 168000 observations on seven disjoint target RX over four dates. This is a fixed-condition mechanism study, not a renewed five-source-count sweep.

Each cell has raw, fixed same-contract G-L and fixed same-contract G-LU readouts. Existing raw and available restoration decisions are reused; nine existing models receive the missing G-L readout. The extension therefore adds 12 F, 21 prediction, 21 independent scoring and one CPU aggregation job, with zero new G. Raw provides the primary F-objective attribution; G-L/G-LU are separately stratified. C=0 plus G-L is a strict L-only system, whereas C=0 plus G-LU is only a supervised F. Variation across three F seeds conditional on a fixed G is not three independently trained complete systems.

Main-table effects are full minus each removed group. The complete eight cells also estimate marginal main effects and AB/AC/BC/ABC interactions within each seed before computing the three-seed mean and sample standard deviation. We retain negative effects and do not infer contributions by adding isolated gains. Metrics include each scene, RX/TX/date, Accuracy, Macro-F1, repaired and introduced errors, and actual segment/total training cost. No new deployment timing is required: A/B/C modify training only and leave the inference graph unchanged. Previously scheduled structural and restoration timing remains separate. The 21 earlier fine ablations are retained as supplementary diagnostics, without using their observed gains to redefine these groups.
